\section{SPN, DES, AES}
\s{SPN}\\
SPN je ponavljanje določenega števila krogov (\(N\)) zaporedja
\[
    \text{XOR} \to S\text{-box} \to P\text{-box}.
\]
\begin{enumerate}
    \item Prištejemo ključ \(K_i\);
    \item Razdelimo besedilo na zloge dolžine \(l\);
    \item Vsak zlog prevedemo v desetiško število in uporabimo \(\pi_S\);
    \item Permutiramo biti s \(\pi_P\).
\end{enumerate}
V zadnjem krogu ponavadi ne naredimo permutacije, saj jo lahko vključimo v ključ \(K_N\).\\
\s{DES}
\begin{itemize}
    \item Velja: \((m_1 + k) - (m_2 + k) \mod n = m_1 - m_2 \mod n\).
    \item DES uporablja \(56\)-bitni ključ: \(c = \text{DES(k, m)}\).
\end{itemize}
\s{AES. Obseg \(\text{GF}(2^8) = \langle x^8 + x^4 + x^3 + x + 1 \rangle\)}\\
Vemo, da se elementi \(\text{GF}(2^8)\) lahko zapišemo kot polinom 
\[
    b_7x^7 + b_6x^6 + b_5x^5 +b_4x^4 +b_3x^3 +b_2x^2 +b_1x + b_0,
\]
kjer so \(b_i \in \Z_2\). Elementi \(\text{GF}(2^8)\) lahko identificiramo z zlogi od 8 bitov. Zlog zapišemo kot par šestnajstiških številk. Na primer
\[
1A = 00011010 = x^4 + x^3 + x
\]
\begin{itemize}
    \item Inverz polinoma poiščemo z razširjenim EA.
    \item Operacije z \(b = 16\): Seštevamo po bitih, množimo polinomi.
\end{itemize}
